\chapter{Nhẹ mà sâu}
\markboth{Nhẹ mà sâu}{Light but Deep}


\section{Chép bài bạn cũng là học — học cách chép không để thầy thấy.}
\begin{truyen}{Chép bài bạn cũng là học — học cách chép không để thầy thấy.}{Humorous but deep.}
\chuhoa{C}{âu thầy nói xong, cả lớp cười — nhưng cười xong nhiều em chợt nghĩ.}
\textit{(The teacher said it; the class laughed — but after the laugh, many suddenly thought.)}
\end{truyen}

\begin{baihoc}
Cười xong nhớ mới là hay.
\textit{(Remembering after the laugh is what counts.)}
\end{baihoc}

\begin{ghinhoanh}
\textbf{Short English to remember:}

Laugh first, then think.

\end{ghinhoanh}

\begin{tuhocnhanh}
1. Câu này khiến em cười vì sao? \textit{Why does this make you laugh?}

2. Sau khi cười em nghĩ ra điều gì? \textit{What do you think after laughing?}
\end{tuhocnhanh}
\ngancach


\section{Thầy biết các em mệt — nhưng mệt mà vẫn học mới là chiến binh.}
\begin{truyen}{Thầy biết các em mệt — nhưng mệt mà vẫn học mới là chiến bin...}{Humorous but deep.}
\chuhoa{C}{âu thầy nói xong, cả lớp cười — nhưng cười xong nhiều em chợt nghĩ.}
\textit{(The teacher said it; the class laughed — but after the laugh, many suddenly thought.)}
\end{truyen}

\begin{baihoc}
Cười xong nhớ mới là hay.
\textit{(Remembering after the laugh is what counts.)}
\end{baihoc}

\begin{ghinhoanh}
\textbf{Short English to remember:}

Laugh first, then think.

\end{ghinhoanh}

\begin{tuhocnhanh}
1. Câu này khiến em cười vì sao? \textit{Why does this make you laugh?}

2. Sau khi cười em nghĩ ra điều gì? \textit{What do you think after laughing?}
\end{tuhocnhanh}
\ngancach


\section{Không hiểu thì hỏi — đừng để trong đầu thành bãi rác.}
\begin{truyen}{Không hiểu thì hỏi — đừng để trong đầu thành bãi rác.}{Humorous but deep.}
\chuhoa{C}{âu thầy nói xong, cả lớp cười — nhưng cười xong nhiều em chợt nghĩ.}
\textit{(The teacher said it; the class laughed — but after the laugh, many suddenly thought.)}
\end{truyen}

\begin{baihoc}
Cười xong nhớ mới là hay.
\textit{(Remembering after the laugh is what counts.)}
\end{baihoc}

\begin{ghinhoanh}
\textbf{Short English to remember:}

Laugh first, then think.

\end{ghinhoanh}

\begin{tuhocnhanh}
1. Câu này khiến em cười vì sao? \textit{Why does this make you laugh?}

2. Sau khi cười em nghĩ ra điều gì? \textit{What do you think after laughing?}
\end{tuhocnhanh}
\ngancach


\section{Em nói đúng rồi — chỉ sai môn thôi.}
\begin{truyen}{Em nói đúng rồi — chỉ sai môn thôi.}{Humorous but deep.}
\chuhoa{C}{âu thầy nói xong, cả lớp cười — nhưng cười xong nhiều em chợt nghĩ.}
\textit{(The teacher said it; the class laughed — but after the laugh, many suddenly thought.)}
\end{truyen}

\begin{baihoc}
Cười xong nhớ mới là hay.
\textit{(Remembering after the laugh is what counts.)}
\end{baihoc}

\begin{ghinhoanh}
\textbf{Short English to remember:}

Laugh first, then think.

\end{ghinhoanh}

\begin{tuhocnhanh}
1. Câu này khiến em cười vì sao? \textit{Why does this make you laugh?}

2. Sau khi cười em nghĩ ra điều gì? \textit{What do you think after laughing?}
\end{tuhocnhanh}
\ngancach


\section{Thầy không la — thầy chỉ đang nói to để em nghe rõ.}
\begin{truyen}{Thầy không la — thầy chỉ đang nói to để em nghe rõ.}{Humorous but deep.}
\chuhoa{C}{âu thầy nói xong, cả lớp cười — nhưng cười xong nhiều em chợt nghĩ.}
\textit{(The teacher said it; the class laughed — but after the laugh, many suddenly thought.)}
\end{truyen}

\begin{baihoc}
Cười xong nhớ mới là hay.
\textit{(Remembering after the laugh is what counts.)}
\end{baihoc}

\begin{ghinhoanh}
\textbf{Short English to remember:}

Laugh first, then think.

\end{ghinhoanh}

\begin{tuhocnhanh}
1. Câu này khiến em cười vì sao? \textit{Why does this make you laugh?}

2. Sau khi cười em nghĩ ra điều gì? \textit{What do you think after laughing?}
\end{tuhocnhanh}
\ngancach


\section{Học một lần chưa đủ — học đến khi não em gật đầu.}
\begin{truyen}{Học một lần chưa đủ — học đến khi não em gật đầu.}{Humorous but deep.}
\chuhoa{C}{âu thầy nói xong, cả lớp cười — nhưng cười xong nhiều em chợt nghĩ.}
\textit{(The teacher said it; the class laughed — but after the laugh, many suddenly thought.)}
\end{truyen}

\begin{baihoc}
Cười xong nhớ mới là hay.
\textit{(Remembering after the laugh is what counts.)}
\end{baihoc}

\begin{ghinhoanh}
\textbf{Short English to remember:}

Laugh first, then think.

\end{ghinhoanh}

\begin{tuhocnhanh}
1. Câu này khiến em cười vì sao? \textit{Why does this make you laugh?}

2. Sau khi cười em nghĩ ra điều gì? \textit{What do you think after laughing?}
\end{tuhocnhanh}
\ngancach


\section{Lớp mình hôm nay im quá — thầy tưởng đi nhầm phòng.}
\begin{truyen}{Lớp mình hôm nay im quá — thầy tưởng đi nhầm phòng.}{Humorous but deep.}
\chuhoa{C}{âu thầy nói xong, cả lớp cười — nhưng cười xong nhiều em chợt nghĩ.}
\textit{(The teacher said it; the class laughed — but after the laugh, many suddenly thought.)}
\end{truyen}

\begin{baihoc}
Cười xong nhớ mới là hay.
\textit{(Remembering after the laugh is what counts.)}
\end{baihoc}

\begin{ghinhoanh}
\textbf{Short English to remember:}

Laugh first, then think.

\end{ghinhoanh}

\begin{tuhocnhanh}
1. Câu này khiến em cười vì sao? \textit{Why does this make you laugh?}

2. Sau khi cười em nghĩ ra điều gì? \textit{What do you think after laughing?}
\end{tuhocnhanh}
\ngancach


\section{Điểm 10 không quan trọng bằng em hiểu được bài.}
\begin{truyen}{Điểm 10 không quan trọng bằng em hiểu được bài.}{Humorous but deep.}
\chuhoa{C}{âu thầy nói xong, cả lớp cười — nhưng cười xong nhiều em chợt nghĩ.}
\textit{(The teacher said it; the class laughed — but after the laugh, many suddenly thought.)}
\end{truyen}

\begin{baihoc}
Cười xong nhớ mới là hay.
\textit{(Remembering after the laugh is what counts.)}
\end{baihoc}

\begin{ghinhoanh}
\textbf{Short English to remember:}

Laugh first, then think.

\end{ghinhoanh}

\begin{tuhocnhanh}
1. Câu này khiến em cười vì sao? \textit{Why does this make you laugh?}

2. Sau khi cười em nghĩ ra điều gì? \textit{What do you think after laughing?}
\end{tuhocnhanh}
\ngancach


\section{Thầy cũng từng là học sinh — nên thầy biết em đang nghĩ gì.}
\begin{truyen}{Thầy cũng từng là học sinh — nên thầy biết em đang nghĩ gì.}{Humorous but deep.}
\chuhoa{C}{âu thầy nói xong, cả lớp cười — nhưng cười xong nhiều em chợt nghĩ.}
\textit{(The teacher said it; the class laughed — but after the laugh, many suddenly thought.)}
\end{truyen}

\begin{baihoc}
Cười xong nhớ mới là hay.
\textit{(Remembering after the laugh is what counts.)}
\end{baihoc}

\begin{ghinhoanh}
\textbf{Short English to remember:}

Laugh first, then think.

\end{ghinhoanh}

\begin{tuhocnhanh}
1. Câu này khiến em cười vì sao? \textit{Why does this make you laugh?}

2. Sau khi cười em nghĩ ra điều gì? \textit{What do you think after laughing?}
\end{tuhocnhanh}
\ngancach


\section{Bài này làm xong em sẽ thấy mình thông minh hơn — hoặc cần uống thêm sữa.}
\begin{truyen}{Bài này làm xong em sẽ thấy mình thông minh hơn — hoặc cần u...}{Humorous but deep.}
\chuhoa{C}{âu thầy nói xong, cả lớp cười — nhưng cười xong nhiều em chợt nghĩ.}
\textit{(The teacher said it; the class laughed — but after the laugh, many suddenly thought.)}
\end{truyen}

\begin{baihoc}
Cười xong nhớ mới là hay.
\textit{(Remembering after the laugh is what counts.)}
\end{baihoc}

\begin{ghinhoanh}
\textbf{Short English to remember:}

Laugh first, then think.

\end{ghinhoanh}

\begin{tuhocnhanh}
1. Câu này khiến em cười vì sao? \textit{Why does this make you laugh?}

2. Sau khi cười em nghĩ ra điều gì? \textit{What do you think after laughing?}
\end{tuhocnhanh}
\ngancach
